\section*{Kurzfassung}
Das elektronische Informationssystem (EIS) der Firma ProSeS versendet automatisierte Benachrichtigungen über verschiedene Kanäle, wie SMS, E-Mail und eine Smartphone-App. Aufgrund der kontinuierlichen Erweiterung des bestehenden Systems ohne begleitende architektonische Anpassungen sind eine hohe Komponentenkopplung und eine eingeschränkte Wartbarkeit entstanden. Ziel dieser Arbeit war daher die Neuentwicklung des EIS.

In einem iterativen Erhebungsprozess wurden elf funktionale Anforderungen priorisiert und 13 nicht-funktionale Anforderungen auf Basis der Wartbarkeitsdefinition der ISO 25010 mit messbaren Abnahmekriterien festgelegt. Darauf aufbauend wurde eine Client-Server-Architektur entworfen. Das Angular-Frontend dient als Konfigurationsoberfläche zur Regelerstellung und -verwaltung. Das C\#-Backend verarbeitet eingehende Ereignisse und steuert den Versand. Für das Backend wurde eine hexagonale Architektur gewählt, die das Hinzufügen neuer Versandkanäle ermöglicht, ohne dass die bestehende Geschäftslogik verändert werden muss.

Die Evaluation zeigt, dass alle als Muss eingestuften funktionalen Anforderungen erfüllt wurden. Die Wartbarkeitsanalyse mittels statischer Analysewerkzeuge belegt im Vergleich zum Altsystem eine höhere Wartbarkeit anhand der ausgewählten Indikatoren. Insgesamt legen die Ergebnisse eine verbesserte strukturelle Grundlage für die langfristige Weiterentwicklung nahe.